\section{Envelopes for singular radial obstacles}
\label{sec:singular-radial-obstacles}\index{Envelope!singular radial obstacle}\index{Regularity!singular envelope residual}

The radial inverse of \cref{sec:radial-residual-regularity} starts with a density. One can instead perturb the singular reference as an obstacle and let the envelope select its measure. The resulting total profile is a convex minorant. Its affine intervals give open regions of zero density, while its contact with the reference at both ends controls the residual across the pole and the hyperplane at infinity.

Keep $k\geq2$, $0<\alpha<1$, the normalized Fubini--Study form $\kappa$ on $\mathbb P^k$ and the potential $u_\alpha$ of \cref{lem:radial-nonlogarithmic-type}. Write $r=|z|^2$ on the chart $Z_0\ne0$, put $x=\log r$, and define
\[
 A(x)=\log(1+\mathrm e^x)+u_\alpha(x).
\]
Here $u_\alpha(x)$ denotes its radial profile. Throughout this section primes denote derivatives in $x$. For $F\in C_c^2(\mathbb R;\mathbb R)$, extend $F(\log r)$ by zero near the pole $o$ and the hyperplane at infinity, and set
\begin{equation}\label{eq:singular-radial-envelope}
 E_F=\left(\sup\{v\in\PSH(\mathbb P^k,\kappa):
       v\leq u_\alpha+F\ \text{everywhere},\ v\preceq u_\alpha\}\right)^*.
\end{equation}
The obstacle takes the value $-\infty$ at $o$. Since $F$ is bounded, its pointwise constraint already implies ordinary domination for each competitor.

\begin{lemma}\label{lem:singular-radial-profile}
The profile $A$ is smooth and strictly convex, with $0<A'<1$, bounded $A''$ and asymptotic slopes $0,1$. Its tangents tend uniformly to $-\infty$ on each compact interval as their point of tangency tends to $-\infty$, and uniformly to the line $y\mapsto y$ as that point tends to $+\infty$. Moreover, $A(x)-x\searrow0$ at the right end.
\end{lemma}
\begin{proof}
As functions of $x$, the quantities $P$ and $T$ of \cref{sec:radial-residual-regularity} are $P=\mathrm e^x/(1+\mathrm e^x)$ and $T=1-\log P$; put $a=\alpha T^{\alpha-1}$. Since $A$ is the profile of $\Phi_\alpha$ and $\mathrm d/\mathrm dx=r\,\mathrm d/\mathrm dr$, \eqref{eq:radial-residual-baseline} and one more differentiation give
\[
 A'=P+(1-P)a,\qquad
 A''=P(1-P)(1-a)+\alpha(1-\alpha)(1-P)^2T^{\alpha-2}.
\]
Since $T\geq1$ and $0<a\leq\alpha<1$, these formulas give the asserted positivity and bounds. On the left, $A(x)=1-(1-x)^\alpha+o(1)$ and $A-xA'$ has leading term $-(1-\alpha)(-x)^\alpha$. On the right,
$A=x+(1-\alpha)\mathrm e^{-x}+O(\mathrm e^{-2x})$ and
$A'=1-(1-\alpha)\mathrm e^{-x}+O(\mathrm e^{-2x})$.
The tangent at $x$ is $y\mapsto A'(x)y+A(x)-xA'(x)$, so both uniform limits follow. Finally $A'-1<0$ gives the decreasing right-end limit.
\end{proof}

\subsection{The convex minorant and its residual}
Put $g=A+F$, $K_F=\|F\|_\infty$, and let $C_F$ be the supremum of all affine functions below $g$ on $\mathbb R$. Since $A-K_F$ is convex and is the supremum of its supporting lines,
\begin{equation}\label{eq:singular-radial-sandwich}
 A-K_F\leq C_F\leq A+F\leq A+K_F.
\end{equation}
In particular $C_F$ is a finite convex function.

\Needspace{16\baselineskip}
\begin{theorem}\label{thm:singular-radial-envelope}
The envelope \eqref{eq:singular-radial-envelope} is
\[
 E_F=C_F(\log r)-\log(1+r)
\]
on the punctured affine chart, with its canonical extensions at both ends. It has ordinary type $u_\alpha$. The residual $R_F=E_F-u_\alpha$ off $o$ extends to $C^{1,1}(\mathbb P^k)$, is zero near $o$ and is smooth near infinity. The function $C_F$ is affine on each component of $\{C_F<g\}$, and, with $L_2=\sup_{\mathbb R}\max(g'',0)$,
\begin{equation}\label{eq:singular-radial-slope-bound}
 0\leq C_F'(b)-C_F'(a)\leq L_2(b-a)\qquad(a<b).
\end{equation}
\end{theorem}
\begin{proof}
We first identify the full envelope, then examine contact and the two ends. Every affine minorant of $g$ has slope in $[0,1]$, because $g\to-\infty$ on the left and $g=x+o(1)$ on the right; hence every supporting slope of $C_F$ lies in $[0,1]$. The sandwich forces its asymptotic slopes to be $0,1$.

The competitor family is invariant under the unitary group $U(k)$ acting on the affine coordinates. Its regularized supremum is therefore radial. Adding smooth local potentials of $\kappa$ gives a locally upper-bounded family of plurisubharmonic functions, so the envelope is $\kappa$-plurisubharmonic. The competitor $u_\alpha-K_F$ prevents collapse, and upper semicontinuity of the obstacle preserves the everywhere bound after regularization, including at $o$. A radial plurisubharmonic total potential $B(\log r)$ has convex profile $B$: restriction to a complex line gives this from radial subharmonicity. Positivity in a direction orthogonal to $z$ gives $B'\geq0$. Conversely a convex nondecreasing profile has nonnegative tangential and radial Hessian eigenvalues $B'/r$ and $B''/r$, first for smooth profiles and then by convex smoothing on annuli. The envelope's total profile is consequently a convex minorant of $g$, and is at most $C_F$.

The same criterion makes $C_F(\log r)$ plurisubharmonic on the punctured chart. Its upper bound near $o$ and \eqref{eq:singular-radial-sandwich} give a removable extension with value $-\infty$. Since $C_F'\leq1$, the function $C_F-x$ is nonincreasing and, by the sandwich, has a finite right-end limit. The proposed $E_F$ is thus locally bounded above near infinity. Add a local smooth potential of $\kappa$ and apply \cref{lem:qpsh-removable-analytic} across the hyperplane. This gives a global competitor, hence the formula. The sandwich proves ordinary equivalence without subtracting the infinite values at $o$.

On a compact interval in the free set the gap $g-C_F$ is positive. Replacing $C_F$ on a sufficiently short subinterval by its chord preserves convexity at both endpoints and stays below $g$. Unless $C_F$ is locally affine this would give a larger minorant. It follows that it is affine on each free component. At a contact point $y$, domination by $g$ gives
$C'_{F,-}(y)\geq g'(y)\geq C'_{F,+}(y)$, and convexity reverses the outer inequality. Thus $C_F$ is differentiable at contacts, with derivative $g'$, and it is differentiable on free intervals as well. Its derivative is continuous. Let $a<b$. If both points lie in one free interval, then $C_F'(b)=C_F'(a)$. Otherwise, if $a$ lies in a free interval, move it to the right endpoint of that interval, and if $b$ lies in a free interval, move it to the left endpoint; these endpoints are finite, since an unbounded free interval containing one of the points would contain both. The resulting contact points $a'\leq b'$ in $[a,b]$ have the same respective derivatives, and
\[
 0\leq C_F'(b)-C_F'(a)=g'(b')-g'(a')
 \leq L_2(b'-a')\leq L_2(b-a),
\]
which proves \eqref{eq:singular-radial-slope-bound}.

For all sufficiently negative $x$, the tangent to $A$ at $x$ lies below $g$ on $\operatorname{supp}F$, by \cref{lem:singular-radial-profile}, and below $A=g$ elsewhere. Hence $C_F=A$ on a left tail. The residual is zero near $o$.

For the other end, put $L_F=\inf_{\mathbb R}(g(x)-x)\leq0$. If this infimum is attained at $y$, the affine minorant $x+L_F$ touches $g$ there. The bound $C_F'\leq1$ forces $C_F=x+L_F$ for $x\geq y$. A negative infimum is always attained, since $g-x$ tends to $+\infty$ on the left and to zero on the right. With $R=1/r$, the envelope on this tail is
\begin{equation}\label{eq:singular-radial-infinity-tail}
 E_F=L_F-\log(1+R).
\end{equation}
If $L_F=0$ is not attained and $F\ne0$, the continuous function $g(y)-y$ has a positive minimum on $\operatorname{supp}F$. The right tangents of $A$ approach the line $y$ uniformly there, so every sufficiently late tangent lies below $g$ and gives $C_F=A$ at its contact point. Thus $C_F=A$ on a right tail; for $F=0$ this holds everywhere. Both alternatives give a smooth residual near infinity, because $R$ and $u_\alpha$ are smooth there.

On the intervening finite annulus,
\[
 |C_F-A|\leq K_F,\qquad |C_F'-A'|\leq1,
 \qquad |(C_F-A)''|\leq L_2+\|A''\|_\infty
 \quad\text{a.e.}
\]
Composing with the smooth coordinate function $x=\log r$ gives $C^{1,1}$ regularity. The exact end formulas extend it over the whole projective space.
\end{proof}

\begin{corollary}\label{cor:singular-radial-envelope-mass}\index{Monge--Ampere measure@Monge--Amp\`ere measure!singular obstacle envelope}
The non-pluripolar measure of $E_F$ has mass one and satisfies
\[
 \mathrm{MA}_\kappa(E_F)(\{a<\log r<b\})=C_F'(b)^k-C_F'(a)^k
 \qquad(-\infty<a<b<\infty).
\]
It vanishes on every free interval. For any $F,G\in C_c^2(\mathbb R;\mathbb R)$,
\begin{equation}\label{eq:singular-radial-value-stability}
 \|R_F-R_G\|_{C^0(\mathbb P^k)}\leq\|F-G\|_\infty.
\end{equation}
\end{corollary}
\begin{proof}
On an annulus, the radial flux formula \eqref{eq:radial-residual-flux}, followed by convex smoothing and bounded Bedford--Taylor continuity, gives the displayed measure. Continuity of $C_F'$ excludes boundary atoms. The non-pluripolar truncations discard $o$, and smoothness at infinity excludes hyperplane mass. The asymptotic slopes $0,1$ therefore give mass one. An affine profile has zero radial eigenvalue and hence zero top-degree measure. Finally, if $\delta=\|F-G\|_\infty$, subtracting $\delta$ from a competitor for $F$ gives a competitor for $G$, and conversely. The resulting value bound for the envelopes gives \eqref{eq:singular-radial-value-stability} off $o$, and continuity of the residuals extends it across both ends.
\end{proof}

For example, take a nonnegative smooth compactly supported bump $\chi$ with $\chi(x_0)=1$ and $F=-M\chi$, where $M>A(x_0)-x_0$. Then $L_F<0$, and \eqref{eq:singular-radial-infinity-tail} gives an open free region of zero density. The envelope still has its sublogarithmic pole, while its residual is globally $C^{1,1}$.

\subsection{Uniform bounds through moving contact}
For a fixed interval $[a,b]$ and $M>0$, let $\mathcal F_{a,b,M}$ consist of real $C^2$ functions supported in $[a,b]$ with
$\|F\|_\infty+\|F''\|_\infty\leq M$. Take all norms in one fixed finite smooth projective atlas.

\begin{theorem}\label{thm:singular-radial-uniform-stability}\index{Envelope!derivative stability}
There is $K=K(k,\alpha,a,b,M)$ such that
\[
 \|R_F\|_{C^{1,1}}\leq K\qquad(F\in\mathcal F_{a,b,M}).
\]
For every $0<\gamma<1$ and $F,G\in\mathcal F_{a,b,M}$ with $\delta=\|F-G\|_\infty\leq1$,
\[
 \|R_F-R_G\|_{C^{1,\gamma}}
 \leq K_\gamma\delta^{(1-\gamma)/2},
\]
where $K_\gamma$ has the same fixed-parameter dependence. These bounds hold across changes of the infinity alternative in \cref{thm:singular-radial-envelope}.
\end{theorem}
\begin{proof}
The left tangents of $A$ lie below $A-M$ on $[a,b]$ once their contact point is less than a common $a_0<a$. They are minorants of every $A+F$, so all residuals vanish on the same neighborhood of $o$. The sandwich and slope bound give uniform real Hessian bounds on any fixed finite annulus, since $L_2\leq\|A''\|_\infty+M$. It remains to control infinity, where the right contact point need not stay bounded.

Put $B=\max(b,0)+1$. For $x>B$, the obstacle is $A$. At contact $C_F'=A'$. On a free interval with finite right contact endpoint $d$, its constant slope is $A'(d)\geq A'(x)$; an unbounded right free interval has slope one. Thus
\[
 0\leq d_F(x):=1-C_F'(x)\leq1-A'(x)\leq\mathrm e^{-x}.
\]
The locally Lipschitz functions $C_F',A'$ agree on the contact set, so their derivatives agree almost everywhere there. Indeed, a differentiable Lipschitz function has derivative zero at almost every density point of its zero set. On the free set $C_F''=0$. The reference formula consequently gives
\begin{equation}\label{eq:singular-radial-weighted-tail}
 0\leq C_F''(x)\leq A''(x)
 \leq(1+\alpha(1-\alpha))\mathrm e^{-x}
 \quad\text{a.e. on }x>B.
\end{equation}

In a chart $Z_j\ne0$, $j\geq1$, use $\zeta=Z_0/Z_j$ and $y=(Z_i/Z_j)_{i\geq1,i\ne j}$. Then
\[
 R=\frac{|\zeta|^2}{1+\sum|y_i|^2},\qquad x=-\log R.
\]
On fixed relatively compact chart pieces, $|\nabla R|\leq c\sqrt R$ and $|\nabla^2R|\leq c$. Write the envelope as
$\mathcal E_F(R)=C_F(-\log R)+\log R-\log(1+R)$ for $R>0$. Its derivatives are
\[
 (\mathcal E_F)_R=\frac{d_F}{R}-\frac1{1+R},\qquad
 (\mathcal E_F)_{RR}=\frac{C_F''-d_F}{R^2}+\frac1{(1+R)^2}
 \quad\text{a.e.}
\]
The tail bounds give $|(\mathcal E_F)_R|\leq2$ and
$|(\mathcal E_F)_{RR}|\leq c_\alpha/R+1$. The real chain rule therefore bounds its Hessian uniformly, and its gradient tends to zero as $\zeta\to0$. Also $L_F=\lim_{x\to\infty}(C_F(x)-x)\in[-M,0]$, and integrating $d_F\leq\mathrm e^{-x}$ gives
$|C_F(x)-x-L_F|\leq\mathrm e^{-x}$. Hence $\mathcal E_F-L_F=O(R)$, so the extended envelope is differentiable on the hyperplane with zero gradient. The Hessian bound makes the gradient Lipschitz off the hyperplane. Apply it along straight coordinate segments and use the continuous zero gradient at a crossing; on segments contained in the hyperplane the gradient is zero. This proves a uniform $C^{1,1}$ extension. Subtract the fixed smooth $u_\alpha$ near infinity and join this estimate to the common pole neighborhood and a fixed annulus.

For the stability estimate put $Z=R_F-R_G$. By \eqref{eq:singular-radial-value-stability}, $\|Z\|_\infty\leq\delta$, and its gradient has a uniform Lipschitz bound. On a smaller fixed chart, a real directional increment of length $h$ gives
$|\partial_e Z|\leq2\delta/h+Kh$. Taking $h$ of order $\sqrt\delta$, or enlarging the constant when limited by the chart radius, yields $\|\nabla Z\|_\infty\leq c\sqrt\delta$. For points $\xi,\xi'$ at distance $t$,
\[
 |\nabla Z(\xi)-\nabla Z(\xi')|\leq\min(2c\sqrt\delta,2Kt).
\]
Dividing by $t^\gamma$ and splitting at $t=\sqrt\delta$ gives the estimate in a finite atlas. When $\delta=0$, the value bound gives equality. This argument does not assert convergence in $C^{1,1}$.
\end{proof}

\subsection{The full envelope on a Hermitian product}
Let $(Y,h_0)$ be a connected compact Hermitian manifold of dimension $m\geq1$, and on $X=\mathbb P^k\times Y$ put
$H=\operatorname{pr}_1^*\kappa+\operatorname{pr}_2^*h_0$, $N=k+m$ and $V_0=\int_Yh_0^m$. Pull $u_\alpha$ and $F$ back to $X$, and define $E_F^H$ by \eqref{eq:singular-radial-envelope} using the full family $\PSH(X,H)$.

\begin{proposition}\label{prop:singular-radial-full-product}\index{Envelope!full Hermitian product}
One has $E_F^H=\operatorname{pr}_1^*E_F$ and
\[
 \mathrm{MA}_H(E_F^H)=\binom Nk\,\mathrm{MA}_\kappa(E_F)\otimes h_0^m,
 \qquad \mathrm{MA}_H(E_F^H)(X)=\binom Nk V_0.
\]
The pulled-back residual is globally $C^{1,1}$, and $E_F^H$ belongs to
$\mathcal E_{\mathrm{env}}(X,H,u_\alpha)$. On each fixed product the uniform bounds and derivative stability of \cref{thm:singular-radial-uniform-stability} hold for these residuals.
\end{proposition}
\begin{proof}
Every competitor restricts to a projective slice as a $\kappa$-plurisubharmonic function or as $-\infty$ identically. To check the exact lower bound when $H$ is nonclosed, choose a smooth quadratic correction $q$ with $\ddc q\geq H$ on a coordinate polydisk and let
$v_\epsilon=(v+q)*\eta_\epsilon-q$ on a smaller polydisk, where $\eta_\epsilon$ is the rescaling of a fixed smooth nonnegative radial probability kernel. Radial plurisubharmonic convolutions decrease to the canonical $v+q$, and
\[
 \ddc v_\epsilon\geq(\ddc q-H)*\eta_\epsilon-\ddc q
 \geq-H-o(1)\omega_{\mathrm{Eucl}}.
\]
Restrict this smooth inequality to a fixed slice. The corrected restrictions decrease either to $-\infty$ or to a plurisubharmonic function in local $L^1$; in the latter case the background error tends to zero distributionally. Subtracting $q$ gives the exact $\kappa$ bound, and these canonical restrictions patch on the connected slice. The everywhere obstacle constraint now gives $v\leq\operatorname{pr}_1^*E_F$ on every slice. Upper regularization preserves this bound. Conversely every base competitor lifts, since $h_0>0$, proving equality of the full envelopes.

The bounded truncations $\max(E_F^H,-j)$ are pullbacks. By \cref{prop:bounded-products,lem:bounded-products-submersion}, multilinearity and dimension count leave only the term with $k$ projective factors and $m$ factors $h_0$. Restrict to strict superlevel sets and let $j\to\infty$ to obtain the stated non-pluripolar measure and its finite mass. The residual regularity and uniform estimates pull back in a fixed smooth product atlas. Finally $E_F^H\geq u_\alpha-K_F$. At $t=1$ take the stretching witness $v_1=E_F^H$; for $t>1$ take $v_t=u_\alpha-tK_F$, so that
\[
 t^{-1}v_t+(1-t^{-1})u_\alpha=u_\alpha-K_F\leq E_F^H.
\]
Together with ordinary domination this is \cref{def:relative-envelope-class}, including the canonical values at the common pole.
\end{proof}

\begin{example}\label{ex:singular-radial-hopf-fourfold}\index{Hopf manifold!singular obstacle product}
There is a nonclosed positive background on a non-K\"ahler fourfold for which the preceding full envelope formula and uniform residual estimates hold.
\end{example}
\begin{proof}
Take $Y=(\mathbb C^2\setminus\{0\})/\langle\zeta\mapsto\mathrm e\zeta\rangle$. With $s=\log|\zeta|^2$, the form $h_0=\ddc s+\mathrm ds\wedge\dc s$ is invariant under the quotient. Its two summands are positive on the tangential and radial complex directions, respectively, so it is Hermitian. Moreover $\mathrm dh_0=-\mathrm ds\wedge\ddc s\ne0$. The Hopf surface is diffeomorphic to $S^3\times S^1$ and has odd first Betti number, so it is not K\"ahler. A K\"ahler form on $X=\mathbb P^2\times Y$ would restrict to one on each fiber $Y$. Thus $X$ is a non-K\"ahler fourfold with $\mathrm dH\ne0$, and \cref{prop:singular-radial-full-product} applies.
\end{proof}
